\documentclass[11pt]{article}
\usepackage{amsmath,amssymb,amsthm}
\newtheorem{theorem}{Theorem}
\title{Exact Antisymmetry Does Not Propagate Independent Child Words}
\author{}\date{}
\begin{document}\maketitle
\paragraph{Status and scope.}
This is a complete written algebraic audit, accompanied by exact rational
checks. It is not Lean-checked. It concerns the additional requirement
that a parent and both literal partial-fraction components have numerators
antisymmetric about the same external midpoint. The actual HeadComplexity
model and the adaptive Haar decoder do not impose that requirement.
Thus the result does not obstruct either accepted global target, nor
does it obstruct the one-node positive-root theorem.

Let $P^\ast(t)=P(T-t)$. Let $M\geq2$ be even, and let $Q_L,Q_R$ be
monic degree-$M$ polynomials with simple, mutually disjoint real roots
strictly below $T/2$. A decomposition into these two pole sets has
numerators
\[
 \frac{N}{Q_LQ_R}=\frac{N_L}{Q_L}+\frac{N_R}{Q_R},\qquad
 N=N_LQ_R+N_RQ_L,\qquad \deg N_L,\deg N_R<M.                 \tag{1}
\]
This is precisely the partial-fraction split: the component residues are
inherited without any extra sign changes. This qualification matters for
the actual Haar decoder. A descendant Haar score applies a fresh signed
mask to the residues in its support, and its numerator need not be $N_L$
or $N_R$ in (1). Therefore the following results do not exclude simultaneous
antisymmetry of differently masked Haar scores. They exclude a shortcut
that propagates symmetry by literal splitting alone.
Assume additionally that
$N^\ast=-N$, $N_L^\ast=-N_L$, and $N_R^\ast=-N_R$.

\begin{theorem}[Coupled child signs]
Under these hypotheses, for every $t>T/2$ the child numerators either
both vanish or have opposite strict signs. The child rational sums
therefore have opposite signs wherever nonzero on the whole exterior
interval. This holds for every partition of the poles.
\end{theorem}
\begin{proof}
Put $D_i=Q_i^\ast-Q_i$. Taking the reflection of (1) and using the three
antisymmetry identities gives
\[
 N_LD_R+N_RD_L=0.                                            \tag{2}
\]
For every root $p<T/2$ and $t>T/2$,
\[
 (t-p)^2-(T-t-p)^2=(2t-T)(T-2p)>0.
\]
Since $t-p>0$, multiplication over all roots gives
$|Q_i(T-t)|<Q_i(t)$, and therefore $D_i(t)<0$.
Equation (2) now proves the assertion. Both denominators in (1) are
strictly positive at exterior points.
\end{proof}

\begin{theorem}[Dimension and shared exterior zeros]
Let $G$ be the monic greatest common divisor of $D_L,D_R$, and write
$D_i=G U_i$, $g=\deg G$. Then $g$ is odd, $1\leq g\leq M-1$, and
all triples satisfying (1) and the three antisymmetries are parameterized
by
\[
 N_L=S U_L,\qquad N_R=-S U_R,\qquad S^\ast=-S,
 \qquad \deg S\leq g.                                      \tag{3}
\]
The solution space has dimension $(g+1)/2$. Every nonzero solution has
at most $(g-1)/2\leq M/2-1$ common child zeros in $t>T/2$.
\end{theorem}
\begin{proof}
If $s_i$ is the sum of the roots of $Q_i$, the leading coefficient of
$D_i$ in degree $M-1$ is $2s_i-MT<0$. Thus each $D_i$ has degree $M-1$.
It is antisymmetric, hence its midpoint zero has odd multiplicity and
all other roots occur in reflected pairs, with multiplicities. The same
holds for their gcd, so $g$ is odd and $G^\ast=-G$; consequently
$U_i^\ast=U_i$.

The polynomials $U_L,U_R$ are coprime. Equation (2) is equivalent to
$N_LU_R+N_RU_L=0$, which forces (3) for a polynomial $S$. The degree
bounds on the child numerators give $\deg S\leq g$; their antisymmetry
is then equivalent to $S^\ast=-S$. Conversely these conditions suffice,
including the parent's antisymmetry by (2). An odd polynomial about
$T/2$ of degree at most the odd number $g$ has $(g+1)/2$ coefficients.
Neither $G$ nor $U_i$ has an exterior real zero, since $D_i(t)<0$ there.
The exterior child zeros are exactly those of $S$. Its nonmidpoint real
zeros occur in reflected pairs, giving the stated bound.
\end{proof}

\paragraph{Polynomial-preimage branches.}
For the proposed split $Q_L=Q_0-\delta$, $Q_R=Q_0+\delta$, the two
reflection differences coincide. Whenever the shifted polynomials have
the required real pole geometry, (2) gives exactly
\[
 N_R=-N_L,\qquad N=2\delta N_L.
\]
Small shifts of a polynomial with simple roots below $T/2$ do have that
geometry by continuity. Thus this split maximizes the symmetry-preserving
freedom, but it still couples the two child words. In particular, exact
antisymmetry does not supply independent child targets even for this
particularly structured alternative to translated pole pairs.

\paragraph{Different child grids do not restore universal freedom.}
On a shared exterior grid the two words must be complementary. If each
child instead has $M/2$ fixed exterior query points and the two grids are
disjoint, their signs, with one child's signs reversed, are evaluations
of the same $S$ up to a fixed common sign. Prescribing alternating signs
on the merged $M$ ordered points would require at least $M-1$ exterior
zeros. The preceding theorem permits only $M/2-1$. This conclusion
continues to hold if the pole partition is chosen adaptively, provided
all three exact antisymmetries and the same external midpoint are
retained. A successful recursive construction based on these literal components
must relax at least one
of those extra requirements, for example by allowing non-antisymmetric
component numerators or additional residue sign masks. The latter masks
are already present in the Haar decoder and remain outside this audit.
\end{document}
